\subsection{非単位重複根atlasの独立な $b^2$ 結合の全排除}
\label{subsec:cubic-third-root-nonunit-independent}
\begin{theorem}[非単位重複根の最初の全領域排除]
\label{thm:cubic-third-root-nonunit-independent-empty}
$\mu_B=(\omega,\omega,1,\omega^2)$ の22tensor atlasを考える。
重み1の自己対称middleの $b^2$ 列を
\[
 p(U)c+c\,r(U),\qquad
 p=(a_{13},a_{20}),\quad r=(-\omega a_9,-\omega a_{17})
\]
と書く。$p,r$ が独立ならH3切断は存在しない。
一方だけ非零の場合も整域性に反する。
従ってこのスペクトルのH3必要候補は、両結合零又は両非零比例に限る。
\end{theorem}
\begin{proof}
一方だけ非零なら $p(U)c$ 又は $c\,r(U)$ が二次関係となるので、整域性に反する。
以下では独立とする。

$U$ の基底変更で $p=(1,0),r=(0,1)$ とできる。
重み1の語順 $(xc,yc,b^2,cx,cy)$ における自己対称blockを $H_0$ とする。
$\operatorname{rank}H_0=2$、$b^2,b^2$ 成分は零、$b^2$ 列は
$P=(1,0,0,0,1)^t$ なので
\[
 H_0=PV^t+VP^t,\qquad V_{b^2}=1
\]
と書ける。これはrank2対称形式と非零isotropic列に対する通常の分解である。
元中山巡回式により、$Uc$ 主blockを $A$ とすると $cU$ 主blockは $-\omega A$。
両者を比較すると
\[
 V=(0,t,1,-\omega t,0)^t.
\]
従ってこの重みの二関係は
\[
 xc+cy=0,\qquad b^2+t yc-\omega t cx=0.
\]
$t=0$ なら純 $b^2$ 関係となるので、整域性から $t\ne0$。

元22tensorの係数は
\begin{gather*}
 a_{13}=1,\ a_{20}=a_9=0,\ a_{17}=-\omega^2,\quad
 a_{11}=a_{19}=a_{14}=a_7=0,\\
 a_{12}=a_3=t,\quad a_{16}=-\omega^2t,
 \qquad a_{10}=m,\ a_{18}=n,\ a_{21}=T
\end{gather*}
となり、残る $U^3b$ の八係数を
$(a_0,a_1,a_2,a_4,a_5,a_6,a_8,a_{15})$ とする。
重み $\omega$ の行順を $(xb,yb,bx,by,c^2)$、
重み $\omega^2$ の列順を $(x^2,xy,yx,y^2,bc,cb)$ としたmiddleを $M$ と書く。
その最後の二列の二つのminorは
\[
 J=mn-\omega-1,
 \qquad mn.
\]
その差は $\omega^2\ne0$ なので、全 $m,n$ でこの二列への射影はrank2である。
もし $T=0$ なら $c^2$ 行は最後の二成分が零、最初の成分は
$-(\omega+1)t\ne0$ となる。
$M$ のrank2と射影rank2に反するので $T\ne0$。

$t=T=1$ と正規化できる。
$x,y$ を共に $\alpha$ 倍、$c$ を $\gamma$ 倍、$b$ を固定し、topを $\eta$ 倍する。
$\gamma^3=(tT)^{-1},\alpha=(t\gamma)^{-1},\eta=t$ と選べば、
$p,r$ の正規化を保ち、$t,T$ を共に1とする。$m,n$ は変わらない。

まず $J\ne0$ とする。$yb,bx$ 行の最後二列をpivotにし、残る三行を消去する。
$J$ を掛けて分母を払った12式は、上の八係数に関する線型系である。
最初の8式は斉次で、その $8\times8$ 行列式は正確に
\[
 \omega J^4 F L,\qquad
 F=m^3+\omega mn-n^3,\quad
 L=m+(\omega+1)(n+1).
\]
$FL\ne0$ なら八係数は全て零となるが、残る $c^2$ 行の $t=1$ 項は消えない。
従って $F=0$ 又は $L=0$ が必要である。

これらの全曲線は有限の特殊点を除く一般点の検算でなく、
次の多項式Smith変換で全値判定する。
$F=0,(m,n)\ne(0,0)$ は
\[
 m=-\frac{\omega q^2}{q^3-1},\qquad
 n=-\frac{\omega q}{q^3-1},\qquad q\ne0,\ q^3\ne1
\]
と書ける。
$L=0$ は $m=\omega^2(q+1),n=q$ と書ける。
それぞれで12式を八未知数の多項式線型系 $A(q)u=b(q)$ にし、
左・右Smith変換 $U A V=D$ を元行列へ掛け戻す。
$U,V$ の行列式は非零定数である。
$D$ の最後4行は零なので、$Ub$ の最後4成分は全て零でなければならない。
これらの成分の多項式gcdはそれぞれ
\[
 q(q^6-q^3+1),\qquad q^2+q+1.
\]
一方、各曲線の $J$ はそれぞれ
\[
 \frac{\omega^2(q^6-q^3+1)}{(q^3-1)^2},\qquad
 \omega^2(q^2+q+1).
\]
従って $J\ne0$ では両曲線とも不可能であり、残り得るのは節点 $(m,n)=(0,0)$ だけである。
行分母を払う際に可逆化したのは、第一曲線では $q^3-1$ のみである。
Smithの途中pivotで零になる値を捨てていない。

次に $J=0$ とする。$m=q\ne0,n=-\omega^2/q$ と置き、
今度は $xb,by$ 行をpivotにする。
同じ多項式Smith変換で、$Ub$ の最後4成分のgcdは $q$ となる。
行分母は $q$ の冪だけなので、この曲線も全値で不可能である。

最後の節点では元12式を係数体上で新しく解き、八係数は正確に
\[
 (a_0,a_1,a_2,a_4,a_5,a_6,a_8,a_{15})
 =(0,-1,0,0,\omega,0,0,0)
\]
となる。
この元22tensorの次数3,4商を全二次関係から構成すると次元は20,35。
中山固定三次7個の正常行列は $36\times7$ であり、rank7の非零minorを持つ。
従って正常核は零で、この最後の点もH3を与えない。

公開再現器は \texttt{verify\_cubic\_third\_root\_nonunit\_independent.py}。
元22tensor、rank2標準形、斉次行列式、三曲線の全Smith恒等式とgcd、
節点の元正常式を全て標数0で再生する。
非単位重複根atlasの比例結合と零結合は、この定理では未解決である。
\end{proof}
